\chapter{Contexte du Projet}

\section{Introduction}

Ce chapitre présente le cadre général dans lequel s'inscrit notre stage de fin d'études. Nous y décrivons l'organisme d'accueil, le contexte institutionnel, la problématique identifiée et les objectifs du projet.

\section{Présentation de l'organisme d'accueil}

\subsection{L'ENSA Béni Mellal}

L'École Nationale des Sciences Appliquées de Béni Mellal (ENSA BM) est un établissement public d'enseignement supérieur rattaché à l'Université Sultan Moulay Slimane (USMS). Créée en 2003 par décret ministériel n° 2-03-200, elle fait partie du réseau des douze ENSA marocaines et forme des ingénieurs d'État dans les domaines des sciences appliquées.

\begin{table}[htbp]
\centering
\begin{tabular}{|l|l|}
\hline
Rubrique & Information \\
\hline
Dénomination & École Nationale des Sciences Appliquées de Béni Mellal \\
\hline
Sigle & ENSA BM \\
\hline
Université & Université Sultan Moulay Slimane (USMS) \\
\hline
Statut & Établissement public d'enseignement supérieur \\
\hline
Création & 2003 — Décret n° 2-03-200 \\
\hline
Effectif étudiants & ~2 500 (2024-2025) \\
\hline
Enseignants & ~120 enseignants-chercheurs \\
\hline
\end{tabular}
\caption{Tableau 2}
\label{tab:table_2}
\end{table}

Table 1:  Fiche signalétique de l'ENSA BM

\subsection{Missions et offre de formation}

L'ENSA BM a pour mission de former des ingénieurs d'État polyvalents, de développer la recherche scientifique et l'innovation technologique, et de contribuer au développement socio-économique de la région Béni Mellal-Khénifra.

L'école propose quatre filières d'ingénieur accessibles aux niveaux Bac+2 et Bac+3 :

\begin{table}[htbp]
\centering
\begin{tabular}{|l|l|l|}
\hline
Code & Filière & Places \\
\hline
TDI & Télécommunications et Développement Informatique & 30 \\
\hline
IACS & Informatique Appliquée et Cybersécurité & 30 \\
\hline
IAA & Ingénierie Agro-Alimentaire & 30 \\
\hline
G2ER & Génie de l'Énergie et Énergies Renouvelables & 30 \\
\hline
\end{tabular}
\caption{Tableau 3}
\label{tab:table_3}
\end{table}

Table 2: Filières d'ingénieur de l'ENSA BM

\subsection{Service de la Scolarité}

Le service de la Scolarité gère l'ensemble des processus administratifs liés aux étudiants : inscriptions, examens, stages, diplomation et présélection des candidats au cycle ingénieur. C'est le principal bénéficiaire et commanditaire de notre plateforme.

\section{Contexte et problématique}

\subsection{Le processus actuel de présélection}

Le processus actuel de présélection à l'ENSA BM se déroule en plusieurs étapes manuelles : (1) réception des dossiers physiques ou par email, (2) vérification manuelle de la conformité des pièces, (3) saisie des notes dans un fichier Excel, (4) tri et classement manuel, (5) convocation à l'épreuve écrite, et (6) saisie des notes d'épreuve et classement final. Ce processus mobilise plusieurs agents pendant plusieurs semaines.

\subsection{Problèmes identifiés}

\begin{itemize}
\item Traitement lent et chronophage : plusieurs semaines pour quelques centaines de dossiers
\item Risque élevé d'erreurs de saisie lors de la transcription manuelle des notes
\item Absence totale de traçabilité des décisions prises
\item Manque de transparence : les candidats n'ont aucune visibilité sur l'état de leur dossier
\item Doublons non détectés : difficulté d'identifier les candidatures multiples
\item Classement approximatif : calculs manuels sujets à erreur
\end{itemize}

\subsection{Le système LMD marocain}

Le système LMD marocain structure les études supérieures en trois cycles. Chaque semestre est évalué sur 20 points. Les mentions sont attribuées selon les seuils suivants :

\begin{table}[htbp]
\centering
\begin{tabular}{|l|l|}
\hline
Moyenne & Mention \\
\hline
$\geq$ 16.00 & Très Bien \\
\hline
$\geq$ 14.00 & Bien \\
\hline
$\geq$ 12.00 & Assez Bien \\
\hline
$\geq$ 10.00 & Passable \\
\hline
\end{tabular}
\caption{Tableau 4}
\label{tab:table_4}
\end{table}

Notre plateforme doit impérativement respecter ces seuils pour le calcul automatique des mentions et du scoring académique.

\section{Objectifs du projet}

\subsection{Objectif principal}

Concevoir et développer une plateforme web complète permettant d'automatiser et de fiabiliser le processus de présélection des candidats aux filières d'ingénieur de l'ENSA BM.

\subsection{Objectifs spécifiques}

\begin{itemize}
\item Automatisation du rejet des dossiers non conformes via un moteur de règles configurable
\item Scoring académique pondéré par filière avec fuzzy matching des matières
\item Classement automatique des candidats intégrant dossier (40\%) et épreuve écrite (60\%)
\item Notifications en temps réel via WebSocket et email SMTP
\item Traçabilité et audit complet du processus de présélection
\item Import des notes d'épreuve écrite au format Excel
\end{itemize}

\section{Périmètre fonctionnel}

\begin{table}[htbp]
\centering
\begin{tabular}{|l|l|}
\hline
Fonctionnalité & Inclus \\
\hline
Inscription et authentification JWT & Oui \\
\hline
Dépôt de dossier multi-étapes & Oui \\
\hline
Moteur de règles de rejet automatique & Oui \\
\hline
Scoring pondéré par filière & Oui \\
\hline
Dashboard responsable avec KPIs & Oui \\
\hline
Import notes Excel & Oui \\
\hline
Notifications WebSocket + SMTP & Oui \\
\hline
Export résultats Excel & Oui \\
\hline
Application mobile native & Non (hors scope) \\
\hline
Intégration Massar & Non (hors scope) \\
\hline
\end{tabular}
\caption{Tableau 5}
\label{tab:table_5}
\end{table}

Table 3:  Périmètre fonctionnel du projet

\section{Méthodologie de travail}

Nous avons adopté une approche agile itérative, organisée en sprints hebdomadaires sur 4 semaines. Chaque sprint produit un incrément fonctionnel testé et validé.

\begin{table}[htbp]
\centering
\begin{tabular}{|l|l|l|}
\hline
Semaine & Phase & Livrables \\
\hline
S1 & Fondations Backend & Modèles Django, API REST, Auth JWT \\
\hline
S2 & Intelligence métier & Moteur de règles, Scoring, OCR \\
\hline
S3 & Frontend React & Pages candidat, Dashboard, Formulaires \\
\hline
S4 & Tests et finalisation & Tests, Corrections, Documentation \\
\hline
\end{tabular}
\caption{Tableau 6}
\label{tab:table_6}
\end{table}

Table 4: Planning prévisionnel du projet

\section{Conclusion}

Ce chapitre a permis de situer le cadre institutionnel de notre stage et de définir clairement la problématique et les objectifs du projet. Le chapitre suivant présentera l'architecture technique et le modèle de données de la plateforme.

